% !TEX TS-program = pdflatex
% !TEX encoding = UTF-8 Unicode

% This is a simple template for a LaTeX document using the "article" class.
% See "book", "report", "letter" for other types of document.

\documentclass[11pt]{article} % use larger type; default would be 10pt

\usepackage[utf8]{inputenc} % set input encoding (not needed with XeLaTeX)

%%% Examples of Article customizations
% These packages are optional, depending whether you want the features they provide.
% See the LaTeX Companion or other references for full information.

%%% PAGE DIMENSIONS
\usepackage{geometry} % to change the page dimensions
\geometry{a4paper} % or letterpaper (US) or a5paper or....
% \geometry{margins=2in} % for example, change the margins to 2 inches all round
% \geometry{landscape} % set up the page for landscape
%   read geometry.pdf for detailed page layout information

\usepackage{graphicx} % support the \includegraphics command and options
\usepackage{amsmath}
% \usepackage[parfill]{parskip} % Activate to begin paragraphs with an empty line rather than an indent

%%% PACKAGES
\usepackage{booktabs} % for much better looking tables
\usepackage{array} % for better arrays (eg matrices) in maths
\usepackage{paralist} % very flexible & customisable lists (eg. enumerate/itemize, etc.)
\usepackage{verbatim} % adds environment for commenting out blocks of text & for better verbatim
\usepackage{subfig} % make it possible to include more than one captioned figure/table in a single float
% These packages are all incorporated in the memoir class to one degree or another...

%%% HEADERS & FOOTERS
\usepackage{fancyhdr} % This should be set AFTER setting up the page geometry
\pagestyle{fancy} % options: empty , plain , fancy
\renewcommand{\headrulewidth}{0pt} % customise the layout...
\lhead{}\chead{}\rhead{}
\lfoot{}\cfoot{\thepage}\rfoot{}

%%% SECTION TITLE APPEARANCE
\usepackage{sectsty}
\allsectionsfont{\sffamily\mdseries\upshape} % (See the fntguide.pdf for font help)
% (This matches ConTeXt defaults)

%%% ToC (table of contents) APPEARANCE
\usepackage[nottoc,notlof,notlot]{tocbibind} % Put the bibliography in the ToC
\usepackage[titles,subfigure]{tocloft} % Alter the style of the Table of Contents
\renewcommand{\cftsecfont}{\rmfamily\mdseries\upshape}
\renewcommand{\cftsecpagefont}{\rmfamily\mdseries\upshape} % No bold!

%%% END Article customizations

%%% The "real" document content comes below...

\title{Elasticity}
%\author{Giuseppe Romano}
%\date{} % Activate to display a given date or no date (if empty),
         % otherwise the current date is printed 
\begin{document}
\maketitle

\section{Getting started}

Elasticity is a Finite Elements solver for mechanical equilibrium problems. It brings features typically developed for structural stability solvers into device modeling. The coupled treatment of the electro-mechanical problem within a unique framework results very useful to explore for multidisciplinary ideas.
In the following we will compute the strain induced by the lattice mismatch in a system comprising a layer of GaN between two contacts of Al$_x$Ga$_{1-x}$N. Once the mesh is created with gmsh (distributed along TiberCAD) we may start to build the input file. First, we have to insert the region section. 
\begin{verbatim}
Device
{
    dimension = 3 
    meshfile = mesh.msh
    mesh_units = 1e-9
    material = AlGaN
    x = 0.2
    x-growth-direction = (-1,0,1,0) 
    y-growth-direction = (-1,2,-1,0) 
    z-growth-direction = (0,0,0,1)
    Region Well  {material = GaN}
}
\end{verbatim}
All the options included in this section, such ad the material, axis growth and alloy concentration, apply for the whole structure.  If we want to specify a region with different properties, we may use the keyword Region and refer to a region name, as specified in the mesh file. The second part of the input file is devoted to the module declaration

\begin{verbatim}
Module elasticity 
{
Physics{
  BodyForce lattice_mismatch
  {
    x-growth-direction = (-1,0,1,0) 
    y-growth-direction = (-1,2,-1,0) 
    z-growth-direction = (0,0,0,1)
    reference_material = AlGaN
    x = 0.2 
  }
   Contact Base{type = clamp}
  }
}
\end{verbatim}
In physics we declare the physical models to be applied to our device. In our case, with BodyForce lattice mismatch we are adding a body force into the system, induced by the lattice mismatch with the reference lattice which in this case is Al0.2Ga0.8N. In order to avoid a free standing device we may want to freeze a surface. With Boundary Clamp we fix, in this case, the region Base.
The third part is simply Simulation{solve = elasticity} where the default name elasticity is used. By default, files for 3D system are written in vtu which can be read from Paraview. Output data about strain are shown below.

\begin{figure}[h!]
\begin{center}
\includegraphics[width=8.0 cm]{strain.png}
%\caption{In plane strain.}
\label{default}
\end{center}
\end{figure}
\begin{figure}[h!]
\begin{center}
\includegraphics[width=10.0 cm]{profile.png}
%\caption{Strain profile.}
\label{default}
\end{center}
\end{figure}
\section{Example: piezoelectric nanogenerator}

In this example we will show how to perform electro-mechanical simulation of a piezoelectric nanogenerators [Wang et al. Science (2006)]. 
The simulation domain consists of a Zinc oxide nanowire (length  = 200 nm, diameter = 50 nm) growth on a ZnO substrate and surroinded by air. The performance of such a device can be obtained by computed the piezopotential produced when the nanowire is subjected to a shear force. 

With the block device we specify the materials and properties associated with the physical regions. In particular, we have the \emph{Nanowire}, \emph{Substrate} and \emph{Air}. 

\begin{verbatim}
Device
{
    dimension = 3 
    meshfile = mesh.msh
    mesh_units = 1e-9
    material = ZnO
    x-growth-direction = (-1,0,1,0) 
    y-growth-direction = (-1,2,-1,0) 
    z-growth-direction = (0,0,0,1)
    
    Region Nanowire  {
    doping_type = donors
    doping_level = 0.035
    }
    
    Region Substrate  {
    doping_type = donors
    doping_level = 0.035
    }
    
    Region Air  {material = Air}
    
}
\end{verbatim}
In this example we use a isotropic stiffness constant and, therefore, a Young modulus and Poisson's ration must be defined. The base of the substrate is clamped whereas a shear force is applied on the top of the nanowire.

\begin{verbatim}
Module elasticity 
{

  Physics
  {
  
   Stiffness isotropic{Young = 129 Poisson = 0.349}
   
   Contact Base{type = clamp}
   
   Contact Top{type = surface_force force = (0.0, 0.0625,0.0}
   
   }
}
\end{verbatim}
The driftdiffusion module computes the piezopotential due to the applied force.  The piezopolarization uses the strain computed by the elasticity module. The base of the substrate is grounded.

\begin{verbatim}
Module driftdiffusion
{
 
  Physics 
  {
    strain_simulation = elasticity 
  
    Contact Base{Voltage = 0.0}

    Polarization piezo{}
   }
   
}
\end{verbatim}
In the simulations block we specify the simulations to be solved. The default simulation name is the name of the module itself. 

\begin{verbatim}
Simulations{solve (elasticity,driftdiffusion)}
\end{verbatim}
Consistently with results published in [Gao et al NanoLetter (2009)] the free carriers tend to screen the output potential. The potential profile is shown below.

\begin{figure}[h!]
\begin{center}
\includegraphics[width=10.0 cm]{nano.png}
%\caption{Strain profile.}
\label{default}
\end{center}
\end{figure}

\section{Theory}
Details about the theory of continuous elasticity applied to device mechanical deformation can be found in [Povolotskiy et al. JAP (2006)]
TiberCAD computes the mechanical deformation of a body subjected to external forces by means of the equilibrium mechanical equation, i.e.
\begin{equation}
 \frac{\partial}{\partial x_j}\sigma_{ij} = f_i
\end{equation}
where $\sigma$ is the stress second-rank tensor and f is the external body force applied to the system. As we will see below, any strain and stress source can be appropriately mapped into an equivalent body force. The stress tensor $\epsilon$ is related to the mechanical strain by means of the constitutive relationship $\sigma_{ij}=C_{ijlk}\epsilon_{lk}$ where $C$ is the stiffness fourth-rank tensor. The strain is related to the displacement u by the expression
\begin{equation}
\epsilon_{lk}=\frac{1}{2}\left(\frac{\partial u_l}{x_k}-\frac{\partial u_k}{\partial u_l} \right )
\end{equation}
Relying on the symmetry of the strain and stiffness tensor the final equation reads as 
\begin{equation}
 \frac{\partial}{\partial x_j}C_{ijlk}\frac{\partial u_l}{\partial x_k} = f_i
\end{equation}
The non-linear strain is computed by applying the mechanical equilibrium equation on the deformed mesh.  The number of the shape iteration is set by the keyword $shape\_iteration$ (default = 0, i.e. no deformation is performed) whereas the maximum tolerated error can be indicated with $shape\_error$ (default = 1e-3).

\subsection{Stiffness constant}
By default the stiffness constant is anisotropic. For the zincoblend crystal structure the independent coefficients are (with the Voigt notation) $C_{11}, C_{12}, C_{44}$:.
\begin{equation}
C = \left(
 \begin{array}{cccccc}
C_{11} & C_{12} & C_{12} & 0 & 0 & 0\\
C_{12} & C_{11} & C_{12} & 0 & 0 & 0\\
C_{12} & C_{12} & C_{11} & 0 & 0 & 0\\
0 & 0 & 0 & C_{44} & 0 & 0\\
0 & 0 & 0& 0 & C_{44} & 0\\
 0 &  0  &  0 & 0 & 0 & c_{44} 
 \end{array}
 \right)
\end{equation}
For the wurtzite structure we have $C_{11}, C_{12}, C_{13}, C_{44}$:
\begin{equation}
C = \left(
 \begin{array}{cccccc}
C_{11} & C_{12} & C_{13} & 0 & 0 & 0\\
C_{12} & C_{11} & C_{12} & 0 & 0 & 0\\
C_{13} & C_{12} & C_{11} & 0 & 0 & 0\\
0 & 0 & 0 & C_{44} & 0 & 0\\
0 & 0 & 0& 0 & C_{44} & 0\\
 0 &  0  &  0 & 0 & 0 & \frac{C_{11} - C_{22}}{2}   
 \end{array}
 \right)
\end{equation}
An isotropic stiffness can be included with the keyword Stiffness isotropic. In this case, the only independent parameters are the Young modulus (Young in GPa) and the Poisson’s ratio (Poisson).
\begin{equation}
C =
\frac{E}{(1+\nu)(1-2\nu)} \left(
 \begin{array}{cccccc}
1-\nu & \nu & \nu & 0 & 0 & 0\\
\nu & 1-\nu & \nu & 0 & 0 & 0\\
\nu & \nu & 1-\nu & 0 & 0 & 0\\
0 & 0 & 0 & \frac{1-2\nu}{2} & 0 & 0\\
0 & 0 & 0& 0 & \frac{1-2\nu}{2} & 0\\
 0 &  0  &  0 & 0 & 0 & \frac{1-2\nu}{2}   
 \end{array}
 \right)
\end{equation}
While the anisotropic model is included by default, the isotropic one must be esplicitally indicated  
\begin{verbatim}
Stiffness isotropic
{
   Young = 129
   Poisson = 0.349
}
\end{verbatim}
 \section{Boundary conditions}
Surfaces forces $\mathbf{f^0}$ are applied by imposing $\sigma_{ij}  n_j= f_i^0$ along the surface with normal $\mathbf{n}$. This boundary condition can be used with the keyword $surface\_force$. The force can be specified in GPa with force. An example is shown below

\begin{verbatim}
Contact base
{
type = surface_force
force = (0,0,0.5) 
}
\end{verbatim}

On the other hand, one may want to fix some surface of the device. The keyword clamp freezes all nodes of a given surface. 
Example:
\begin{verbatim}Contact substrate {type = clamp}\end{verbatim}.

\section{Body forces}

A constant value body force can be included by means of the keyword \emph{Body\_force constant}.
Example: \begin{verbatim}BodyForce constant {force = (0,1,0)}\end{verbatim}.
When two crystals with different crystal structure are put in contact the lattice mismatch between them may induce a strain $\epsilon^LM$ and, therefore, a stress $\sigma =C\epsilon^{LM}$. This stress contribution acts as a body force 

\begin{equation}
 f_i = -\frac{\partial }{\partial x_j} C_{ijlk}\,\epsilon_{lk}^{LM}
\end{equation}
The strain source can be computed only once the reference lattice is identified. The reference material and its growth axis can be included following the same syntax of the device section.

Example:

\begin{verbatim}
BodyForce lattice_mismatch
    {
     reference_material = AlGaN
     structure = wz
     x = 0.2
     x-growth-direction = (1,0,0,0)
     y-growth-direction = (0,1,-1,0)
     z-growth-direction = (0,0,0,1)
 }
\end{verbatim}
In presence of an electric field there might develop an additional stress source due to the so-called converse piezoelectric effect, given by
\begin{equation}
 \sigma_{il}^{SC}=-e_{ijk}E_k
\end{equation}
The converse piezoelectric effect can be included with the keyword BodyForce piezo and an electrostatic simulation must be indicated. The effective body force reads as
\begin{equation}
 f_i = -\frac{\partial}{\partial x_j} \sigma_{ij}^{CP}
\end{equation}
Example: \begin{verbatim}BodyForce converse_piezo{poisson_simulation = dd} \end{verbatim}. 
Considering all the abovementioned force sources, the plotted total stress and strain are
\begin{equation}
 \sigma_{ij} = C_{ijlk}\left( \frac{\partial u_l}{\partial x_k} + \epsilon_{lk}^{LM} \right) -e_{ijk}E_k  
\end{equation}
\begin{equation}
\epsilon_{kl}=\frac{1}{2}\left(\frac{\partial u_l}{x_k}-\frac{\partial u_k}{\partial u_l} \right ) + \epsilon_{lk}^{LM}
\end{equation}

\newpage
\section{Reference Guide}


\begin{verbatim}
Stiffness/Isotropic
Kind: bulk 
Name: Stiffness
Type: isotropic
database_section: none
Options: 
   young  (double,0.0)
   poisson (double [-0.5,1.0],0.0)

Stiffness/Anisotropic
Kind: bulk 
Name: Stiffness
Type: anisotropic
database_section = stiffness 
Options:
   none

Body Force/constant
Kind: bulk 
Name: BodyForce
Type: constant
database_section: none
Options: 
   force (double vector,(0.0,0.0,0.0))

Body Force/Lattice Mismatch
Kind: bulk 
Name: BodyForce
Type: lattice_mismatch
database_section: none
Options:
     x (double [0.0,1,0],0.0)
     x-growth-direction (double vector)
     y-growth-direction (double vector)
     z-growth-direction (double vector)

Body Force/Converse
Kind: bulk 
Name: BodyForce
Type: converse
database_section: piezoelectricity
Options: 
   poisson_simulation (string,”none")

Boundary/clamp
Kind: interface
Name: Contact
Type: converse
database_section: none
Options: 
   poisson_simulation (string,”none")

Boundary/Surface Force
Kind: interface
Name: Contact
Type: surface_force
database_section: none
Options: 
   force (double vector,(0.0,0.0,0.0)

\end{verbatim}




\end{document}
